\chapter{总结与展望}

本文以全球水稻种植区动态遥感监测及其地表降温效应评估为核心，围绕水稻生长季事件识别、全球水稻时空格局刻画以及稻田地表温度效应量化三个方面展开，
形成了由方法开发、数据产品构建、时空动态分析到地表气候效应评估的完整研究链条。
本章在前文研究基础上，归纳本文主要研究结论，凝练方法、数据和科学应用层面的创新点，并进一步提出全球稻田动态遥感监测与稻田气候效应研究的未来发展方向。

\section{主要研究结论}

第一，本文构建了面向全球复杂稻作制度的 MPD\_DTW 水稻生长季事件识别框架。
该框架将移动物候检测与动态时间规整曲线匹配相结合，利用水稻移栽或返青初期的淹水信号、生长季内植被指数曲线形态以及收获后的绿度衰退过程，逐像元识别水稻生长季起始日和结束日。
与固定时间窗口或单一阈值方法相比，MPD\_DTW 能够在统一方法结构下处理不同地区播栽期漂移、生长季长度差异和年内多季种植并存等问题，使水稻制图对象由静态土地覆盖类别扩展为具有明确起止时间的作物生长季事件。
该方法为全球尺度长期水稻动态监测提供了具有物候解释性和跨区域适应性的技术基础。

第二，本文基于 MPD\_DTW 框架构建了 2001--2022 年全球 500 m 日尺度水稻数据集 GlobalRice500。
该数据集同时提供年度水稻分布、逐像元 SOS、EOS 和年内种植季信息，使全球水稻研究能够从年度面积表达推进到生长季过程表达。
精度验证结果表明，GlobalRice500 在全球典型稻区具有较好的空间识别能力，并与 FAO 国家尺度统计面积保持较高一致性。
基于该数据集，本文揭示了 21 世纪以来全球水稻收获面积总体呈波动增加趋势，至 2022 年净增加超过 2100 万公顷；亚洲长期贡献全球水稻面积的 85\% 以上，非洲则表现出最快的扩张速度。
这些结果表明，GlobalRice500 能够支撑全球、洲际和国家尺度水稻时空动态分析，也为后续稻田环境效应评估提供了统一的空间和时间边界。

第三，本文在 GlobalRice500 提供的真实水稻生长季约束下系统量化了全球稻田地表降温效应，并揭示了其生长季动态和景观结构调控特征。
基于稻田与邻近非水稻农田的对照比较，2003--2020 年全球稻田在水稻生长季日间相对于周边非水稻农田表现出稳定的局地地表降温效应，日间平均 $\Delta LST$ 年际变化范围约为 0.21--0.27 $^{\circ}\mathrm{C}$。
夜间 $\Delta LST$ 接近零，说明稻田降温主要发生在太阳辐射驱动下的日间能量分配过程中。
从生长季内过程看，日间 $\Delta LST$ 呈现明显的阶段性变化，降温峰值主要出现在 DOS 60--100 附近，对应水稻冠层快速发育和蒸散发较强的生长阶段。
该结果表明，稻田降温效应具有明确的物候依赖性，真实生长季约束对于全球尺度稻田地表温度效应评估具有关键作用。

进一步分析表明，稻田地表降温效应受到局地景观结构和空间邻近关系的显著调控。
稻田比例越高的局地景观通常表现出更强的日间降温效应，稻田比例每增加 1\%，日间 $\Delta LST$ 平均增加约 0.0014 $^{\circ}\mathrm{C}$。
随着稻田田块规模由约 0.5 km 增至约 5.5 km，日间平均降温由约 0.13 $^{\circ}\mathrm{C}$ 增至约 0.43 $^{\circ}\mathrm{C}$，生长季峰值降温由约 0.18 $^{\circ}\mathrm{C}$ 增至约 0.69 $^{\circ}\mathrm{C}$。
非水稻农田梯度环带分析进一步表明，稻田降温效应可以向邻近非水稻农田扩展，并随距离增加逐渐减弱。
这些结果共同说明，稻田地表降温并非单个像元的孤立属性，而是具有明确生长季过程、景观尺度调控和空间外溢特征的局地地表热环境效应。


\section{创新点}

本文的创新体现在方法构建、数据产品和科学应用三个相互衔接、层层递进的方面。
MPD\_DTW 通用水稻识别框架解决了全球复杂稻作制度下水稻真实生长季难以统一识别的问题，
GlobalRice500 数据集将该识别能力转化为长期、连续、可复用的全球水稻动态数据产品，
稻田降温效应分析进一步揭示了稻田对地表热环境的生物物理调节作用。
三者共同构成了从作物动态遥感制图到稻田地表气候效应评估的完整研究链条。

第一，本研究提出了面向全球复杂稻作制度的水稻生长季事件识别框架。
全球水稻种植受到气候背景、水分管理、播栽方式、品种熟期和复种制度共同影响，
不同地区水稻生长季的起止时间、持续长度和年内种植季次数差异显著。
传统固定月份窗口、区域作物日历或单一物候阈值方法难以同时适应热带多熟区、
亚热带双季稻区、温带单季稻区以及错季种植区的复杂时序特征。
针对这一问题，本研究提出 MPD\_DTW 框架，将移动物候检测与动态时间规整相结合，
以逐像元、逐日搜索方式识别候选 SOS 和 EOS，
并通过候选生长季曲线与标准水稻生长曲线的形态匹配判断该时段是否符合水稻生长过程特征。
该框架将水稻遥感制图对象由年度静态类别拓展为具有明确起止时间的生长季事件，
使单季稻、双季稻、多季稻和错季水稻能够在统一逻辑下被连续识别。
这一方法提高了全球复杂稻作制度下水稻动态识别的物候解释性和时序适应能力，
也为其他具有明确物候阶段和稳定时序结构的作物动态监测提供了方法参考。

第二，本研究构建了兼具全球覆盖、长期连续、500 m 空间分辨率和逐像元生长季信息的 GlobalRice500 数据集。
已有全球或大区域水稻产品通常难以同时满足全球覆盖、长期序列、中等空间分辨率和高时间分辨率等需求；
部分产品侧重年度面积或静态空间分布，难以表达水稻真实生长季过程。
GlobalRice500 提供 2001--2022 年全球年度水稻分布、逐像元 SOS、EOS 和年内种植季信息，
将水稻空间范围、种植季次和生长季过程统一到同一数据体系中。
与主要表达年度面积或静态空间格局的水稻产品相比，
GlobalRice500 使全球水稻研究由年度面积表达拓展至逐像元生长季过程表达。
该数据集不仅支撑了全球水稻面积变化、洲际贡献和国家尺度格局分析，
也为稻田地表温度效应评估提供了真实生长季约束和一致的空间基础。
因此，GlobalRice500 的贡献不仅在于提供新的全球水稻分布图，
更在于为水稻生产过程、复种制度和稻田环境效应研究提供了可复用的数据底座。

第三，本研究在真实水稻生长季约束下系统揭示了全球稻田地表降温效应及其景观结构调控特征。
已有稻田气候效应研究长期集中于甲烷排放等生物地球化学增温效应，
而稻田通过水层蒸发、冠层蒸腾和能量分配调整产生的地表生物物理效应缺乏全球尺度量化。
基于 GlobalRice500 提供的 SOS--EOS 信息，
本文构建了稻田与邻近非水稻农田之间的 $\Delta LST$ 比较框架，
降低了固定月份窗口在跨气候区、多熟制和错季种植条件下造成的时间错配。
在此基础上，本文量化了全球稻田生长季日间地表降温效应，
揭示了其昼夜差异、生长季动态、气候分区差异、稻田比例调控、田块规模效应和非水稻农田梯度环带所反映的空间外溢特征。
该分析将稻田地表温度效应由单一平均量拓展为具有生长季过程、空间连续性和邻近传播特征的局地地表热环境过程，为理解农业景观结构对地表生物物理效应的调控作用提供了全球尺度证据。
同时，本研究讨论了稻田地表降温效应与甲烷排放升温效应的解释边界，有助于在粮食生产、水资源消耗、温室气体排放和局地热环境调节之间，更全面地认识稻田生态系统的多维气候作用。


\section{研究展望}
前文已分别讨论 MPD\_DTW 方法、GlobalRice500 数据产品以及稻田地表降温效应评估的适用边界。
总体而言，全球水稻遥感监测仍需在空间分辨率提升、生长过程表达、田间管理信息反演和地表气候效应归因等方面进一步深化。
因此，未来研究可围绕多源遥感融合、过程型水稻产品构建、能量平衡机制解析、因果归因框架完善以及气候--环境--粮食综合评估等方向展开。

% 提高空间分辨率
第一，融合多源遥感数据，发展更高空间分辨率和更强云雨适应性的全球水稻动态产品。
GlobalRice500 在全球覆盖、长期连续和逐像元生长季表达方面具有优势，
但 500 m 空间分辨率仍会限制小田块、梯田化稻区和景观破碎区域的精细刻画。
Sentinel-1 SAR 对稻田水分状态和冠层结构变化具有较强敏感性，Sentinel-2 与 Landsat 光学影像能够提供更精细的植被生长信息。
未来可在保持长时间序列一致性的基础上，发展面向不同时间阶段和空间尺度的全球水稻动态产品。
跨传感器辐射一致性、SAR–光学时序融合、云雨条件下观测缺失重建以及不同分辨率产品之间的尺度转换，将是高分辨率全球水稻动态产品构建中的关键问题。

% 物候
第二，在 GlobalRice500 基础上进一步发展过程型水稻遥感产品。
除 SOS 和 EOS 外，未来可进一步提取移栽或返青期、分蘖期、抽穗期、成熟期、收获期、生长季长度、复种强度以及年内种植季衔接关系等变量，使水稻遥感产品从空间分布表达扩展到生长过程表达。
对于稻田环境效应研究而言，田间水分状态和灌溉管理信息尤为关键。若能结合 SAR 后向散射、光学水分指数、地形水文条件和区域农田管理资料，识别连续淹水、间歇灌溉、雨养稻和旱直播稻等主要管理差异，将有助于解释水稻产量形成、甲烷排放、蒸散发变化和 LST 响应之间的耦合关系。
因此，全球水稻遥感监测的下一步重点可由稻田分布与生长季识别进一步拓展为“物候过程—水分状态—环境效应”的综合刻画。

% 全面评估温度效应
第三，综合评估稻田地表降温效应与温室气体增温效应。
本文从地表生物物理过程角度揭示了稻田在水稻生长季的日间降温效应，但稻田气候影响还包括甲烷排放所产生的温室气体增温效应。未来研究可在统一时空框架下，将稻田蒸散发增强和能量分配改变引起的地表降温与甲烷排放引起的辐射增温效应结合起来，全面评估稻田系统对气候的净影响。
这一综合评估框架有助于在保障粮食生产的前提下，权衡稻田节水、减排和热环境调节效应，为全球变化背景下稻田系统可持续管理提供更完整的科学依据。

% 其他环境效应
第四，结合遥感数据与过程模型，综合评估稻田系统的多维环境效应。
本文基于 GlobalRice500 揭示了全球水稻时空动态及其地表降温效应，为认识稻田地表生物物理气候作用提供了数据基础。
未来研究可进一步将 GlobalRice500 及其后续高分辨率产品与作物生长模型、陆面过程模型、水文模型、生物地球化学模型和区域气候模式相结合，评估稻田面积变化、复种制度调整和田间管理差异对粮食生产、水资源利用、蒸散发、甲烷排放、地表温度和近地气候的综合影响。
现有模型对稻田空间分布、物候边界、复种制度、水分管理和能量交换过程的表达仍较为简化；
高时空分辨率水稻动态数据可为模型提供更真实的稻田范围、生长季边界和年内种植季信息，从而改进稻田水文过程、碳氮循环过程和地表能量交换过程的参数化表达。
该方向有助于从粮食供给、资源消耗、温室气体排放和气候反馈等多个维度认识水稻系统，并为全球变化背景下稻田可持续管理提供科学依据。

总体而言，本文证明了全球一致的水稻动态遥感信息对于解释稻田地表生物物理效应的重要基础作用。
未来若能在更高空间分辨率、更完整物候过程、更精细水分状态信息和更强机制约束之间形成协同，全球稻田研究将从水稻空间监测进一步走向粮食生产、资源消耗、生物地球化学排放和地表气候反馈的综合评估；
有助于完善稻田气候效应认识，也可为气候变化背景下农业土地系统的可持续管理提供更系统的科学依据。